\begin{textsample}{POS Dim 9 – global\_south\_2024\_10 – Score 1.00 – global\_south\_2024\_10/115260938.txt}  \label{ex:f9_pos_171}
Israeli defence expert Yossi Kuperwasser said the two-State solution was unrealistic at this time and ceasefire talks in Gaza were delayed.
Israel is concentrating on addressing threats from Hezbollah and Iran, the expert said, highlighting plans to replace Hamas leadership in Gaza and restore Lebanese state control in Lebanon.
ANIA war-ravaged locality Tel Aviv : With the West Asia conflict opening up on multiple fronts, and Israel ' s focus shifting to fighting Hezbollah and Iran, Israeli defence expert, Brig.
Gen ( Res ) Yossi Kuperwasser, believes that ceasefire talks in Gaza are now on the backburner and the ' two-State solution ' is not going to happen \textbf{anytime} soon.
Brig.
Gen. ( Res. ) Kuperwasser is a senior researcher at the Jerusalem Center for Security and Foreign Affairs.
He was the former head of the research division in IDF Military Intelligence and Director General of the Israel Ministry of Strategic Affairs.
In a telephonic interview with ANI, Kupperwasser highlighted that Israel ' s aim is to replace the leadership in both Gaza as well as Lebanon, and also emphasised @ @ @ @ @ @ @ @ @ @ size"for attaining any"realistic solution"in the conflict."A realistic solution is that Hamas and all the Hamas is removed from power in Gaza, that there ' s a new state of affairs in Lebanon, where the Lebanese state goes back into being the owner of the monopoly of the use of force for its territory, which is not the case now,"Kupperwasser said."And when Iran goes back to its real size and start being a power... superpower in the Middle East and focuses on improving the living conditions in Iran, the Iranians, not building its hegemony in the Middle East,"he said.
The former Israeli military intelligence official further stated that Israel plans to"replace"the Hamas leadership in Gaza, and ensure the safety of its citizens in the northern border, in reference to Lebanon."We are planning on replacing the government in Gaza and putting somebody else in place of Hamas.
We plan to make sure that our @ @ @ @ @ @ @ @ @ @ they are back at their homes,"Kupperwasser said."We wish to see the international community, led by the United States, take the necessary steps in order to put back Iran in its real size, not threatening the entire stability of the Middle East.
That ' s what we plan,"the defence expert said..
Hamas launched a massive terror attack against Israel on October 7 last year, killing over 1200 people and holding over 250 as hostages, out of whom around 100 are still in captivity.
In response, Israel launched a strong counter-offensive, targeting Hamas units in the Gaza Strip.
However, the mounting civilian toll has raised concerns over the humanitarian situation in the region.
More than 35,000 Palestinians have been killed in the conflict, according to the Gaza health ministry.
Amid the escalating situation, which many have called a"humanitarian crisis,"almost all big powers including the United States have called for a ceasefire in Gaza and the release of hostages. @ @ @ @ @ @ @ @ @ @ to attain sustainable peace in the Israel-Palestine conflict.
Being asked if there will be a change in Israel ' s strategy if there is a change of guard in the United States with the presidential elections scheduled to be held in November, Kuperwasser said that there will be no change in the US support for Israel."Well, I don't think that they ' re going to stop us from doing what ' s necessary for our security.
And that ' s why we showed the world what Hezbollah was planning to carry out, from what kind of attack it was planning to carry out from southern Lebanon so that people understand that what we do is a necessary precaution in order to prevent a war, prevent a massive terror attack against our population.
So, I don't think they ' re going to stop us from doing that,"the defence expert said.
Referring to the ceasefire talks in Gaza -- which the US and many other powers have been pushing for -- the Israeli defence expert acknowledged that @ @ @ @ @ @ @ @ @ @ can only be thought of after Israel is finished dealing with Hezbollah."They ( US ) would like to reach some sort of a ceasefire down the road, but they understand this is now out of the question, and first of all Israel has to defend itself,"Kupperwasser said."I think it ' s ( ceasefire talks ) in a backburner right now, and maybe after we deal with Iran and after we deal with Hezbollah, we should go back to Hamas, and they will be in a much weaker position.
So, then we should be able maybe to make progress towards an agreement, some sort of a deal that would allow us to get back our hostages from Hamas,"he added.
When asked about the two-State solution, Kupperwasser said that speaking of the issue now is"\textbf{delusional}".
Referring to the October 7 attacks, he said that Palestinians are totally committed to"annihilating the state of Israel.""Well, the two-State @ @ @ @ @ @ @ @ @ @ solution at this point.
And that ' s why it ' s premature totally to speak about it.
The problem is that the Palestinians, as we saw on the October 7, are totally committed to annihilating the state of Israel.
And that ' s what their purpose in life is, to struggle against Israel until its demise.
So, as long as this is the case, speaking about the two-state solution is \textbf{delusional}.
It ' s not going to happen \textbf{anytime} soon,"the Israeli defence expert further said.

% matched lemmas: anytime, delusional
\end{textsample}
